\section{HexCore Phases 56--57: Real-File Projects and Relational World Models}
\label{sec:hexcore-phases-56-57}

\subsection{Programme objective}

Phases 56 and 57 connect AION's open-world ingestion system to persistent
project execution and learned models of change.  Phase 56 establishes that
AION can execute a checksum-bound project across interruptions and revise its
work when a real source file changes.  Phase 57 then replaces a fixed
revision rule with a learned relational theory that predicts which project
artifacts will become stale.

The combined governed chain is:
\[
\begin{aligned}
\text{real files}
&\rightarrow \text{persistent project graph}
\rightarrow \text{checkpointed conclusions}\\
&\rightarrow \text{observed source revision}
\rightarrow \text{learned transition model}\\
&\rightarrow \text{selective invalidation or abstention}
\rightarrow \text{independent verification}
\rightarrow \text{CAU retention}.
\end{aligned}
\]

Neural or symbolic proposals do not become authoritative merely because they
are plausible.  Current source hashes, exact evidence recovery, independent
execution and Cognitive Authority Unification (CAU) remain the acceptance
boundary.

\subsection{Phase 56: real-file long-horizon project execution}
\label{subsec:phase56-real-file-projects}

\subsubsection{Capability contract}

Phase 56 strengthens the earlier synthetic long-horizon benchmark by
operating on real repository result files and technical reports.  Each
project must:
\begin{enumerate}
    \item bind its objective to checksum-addressed source files;
    \item construct and persist an explicit task-dependency graph;
    \item preserve exact evidence locations and knowledge status;
    \item checkpoint after every milestone;
    \item recover after forced process reconstruction;
    \item detect a source revision after producing a draft;
    \item invalidate only work that depends on the changed source;
    \item reconcile the revision or abstain on unresolved contradiction; and
    \item pass a verifier that independently rereads the current files.
\end{enumerate}

\subsubsection{Development and sealed cohorts}

The development cohort used real outputs and reports from Phases 48--52:
external evaluation, document intelligence, project decomposition, tool
invention and continual self-improvement.  The source-disjoint sealed cohort
used Phases 53--55:
\begin{itemize}
    \item blind external evaluation;
    \item open-world ingestion; and
    \item recursive Photon synthesis.
\end{itemize}

Each sealed project contained an authoritative JSON result and its Markdown
technical report.  The repository sources were read-only.  All revision
tests used checksum-bound isolated copies.

\subsubsection{Persistent project representation}

Each entry in \texttt{real\_file\_projects} contains:
\begin{itemize}
    \item project identifier, cohort, family and objective;
    \item milestone plan and explicit dependencies;
    \item source paths, roles and SHA-256 snapshots;
    \item exact report evidence and character offsets;
    \item unresolved questions;
    \item invalidated and repeated tasks;
    \item draft and final manifests;
    \item restart count;
    \item independent-verification record; and
    \item an append-only decision trace.
\end{itemize}

The seven-task project graph is:
\[
\begin{aligned}
\text{result source}
&\rightarrow \text{status extraction},\\
\text{report source}+\text{status}
&\rightarrow \text{evidence binding},\\
\text{status}+\text{evidence}
&\rightarrow \text{manifest},\\
\text{manifest}
&\rightarrow \text{independent verification}
\rightarrow \text{closure}.
\end{aligned}
\]

Every milestone is atomically committed through the persistent HexCore store.
The benchmark reconstructs the runtime after every milestone, forcing six
restart recoveries per sealed project.

\subsubsection{Revision protocol and selective invalidation}

After AION creates and persists its draft manifest, the sealed file controller
amends a copied result file.  The amendment is absent when the draft is
created.  AION must discover the new checksum, supersede stale conclusions
and rebuild the affected artifacts.

A result revision invalidates five downstream tasks:
\begin{enumerate}
    \item status extraction;
    \item evidence binding;
    \item manifest synthesis;
    \item independent verification; and
    \item project closure.
\end{enumerate}

Registration and independent report-ingestion work remain valid.  Relative
to replaying all seven tasks, the bounded saving is:
\[
R_{\mathrm{Phase\,56}}
=
1-\frac{5}{7}
=
0.2857
=
28.57\%.
\]

This measurement establishes dependency-aware revision.  It is not presented
as a general estimate of savings on arbitrary projects.

\subsubsection{Independent disk-rereading verifier}

The final verifier does not trust the saved manifest.  It rereads the current
files and independently checks:
\begin{itemize}
    \item phase identity;
    \item current projected status;
    \item current result and report hashes;
    \item exact recovery of the cited report substring;
    \item recorded source-change detection; and
    \item closure of every required question.
\end{itemize}

The Phase 55 project is an unresolved-conflict negative control.  Correct
behaviour is explicit abstention rather than forced finalization.

\subsubsection{Phase 56 results}

\begin{table}[htbp]
\centering
\begin{tabular}{lr}
\hline
\textbf{Measure} & \textbf{Result}\\
\hline
Sealed projects & 3\\
Real sealed source files & 6\\
Project outcome accuracy & 100\%\\
Weakest-family accuracy & 100\%\\
Source-change detection & 100\%\\
Restart recovery & 100\%\\
Provenance completeness & 100\%\\
Conflict-control abstention & Passed\\
Re-execution reduction & 28.57\%\\
Original repository files modified & 0\\
\hline
\end{tabular}
\caption{Phase 56 sealed real-file project results.}
\label{tab:phase56-results}
\end{table}

CAU promoted:
\[
\texttt{procedure\_real\_file\_projects\_3aff718b5c0d}.
\]

The complete HexCore history passed \(61/61\) tests.  The implementation was
committed as Git revision \texttt{a3c33469} and synchronized with the remote
\texttt{main} branch.

\subsubsection{Phase 56 boundary}

Phase 56 demonstrates persistent real-file project execution, source-change
detection, dependency-aware revision, restart recovery and verified
abstention.  The objective, seven-stage task schema, controlled revision and
final verifier remain engineered.  The result does not establish unrestricted
autonomous project management, arbitrary application control or external
certification.

\subsection{Phase 57: richer learned relational world models}
\label{subsec:phase57-world-models}

\subsubsection{Research objective}

Phase 57 asks whether AION can learn how changes propagate rather than rely
on a hard-coded invalidation rule.  It attaches observed event streams to
real-file projects and requires the learner to:
\begin{itemize}
    \item represent typed source and artifact entities;
    \item compare competing transition explanations;
    \item preserve posterior uncertainty;
    \item predict withheld revision and reconciliation outcomes;
    \item transfer its theory to unseen dependency topologies;
    \item recognise events outside the current model family; and
    \item persist models, criticisms and outcomes across restart.
\end{itemize}

\subsubsection{Relational project state}

The real-file substrate contains JSON results, Markdown reports and YAML
governance policy.  A representative dependency graph is:
\[
\begin{aligned}
\text{result}
&\rightarrow \text{status},\\
\text{report}+\text{status}
&\rightarrow \text{evidence},\\
\text{policy}+\text{status}
&\rightarrow \text{authority check},\\
\text{status}+\text{evidence}+\text{authority}
&\rightarrow \text{manifest},\\
\text{manifest}
&\rightarrow \text{verification}
\rightarrow \text{closure}.
\end{aligned}
\]

Two sealed projects include an additional publication-receipt entity, giving
the learner topologies absent from development.

\subsubsection{Competing model family}

AION compares twelve inspectable world models formed from the Cartesian
product of four invalidation operators and three conflict policies:
\[
\mathcal{M}
=
\{
\text{none},\text{direct},\text{transitive},\text{blanket}
\}
\times
\{
\text{ignore},\text{rebuild},\text{abstain}
\}.
\]

For model \(m\), development scoring combines field accuracy, exact transition
accuracy and a complexity penalty:
\[
\ell(m)
=
8A_{\mathrm{field}}(m)
+4A_{\mathrm{exact}}(m)
-0.03C(m).
\]

Posterior weights are obtained by normalized exponentiation:
\[
p(m\mid D)
=
\frac{\exp(\ell(m)-\ell_{\max})}
{\sum_{m'\in\mathcal{M}}\exp(\ell(m')-\ell_{\max})}.
\]

The selected theory is:
\[
\boxed{
\text{transitive dependency propagation}
+\text{conflict abstention}
}
\]
with posterior probability \(71.28\%\).  AION also retains:
\begin{itemize}
    \item transitive propagation plus rebuild: \(17.37\%\); and
    \item transitive propagation plus ignore: \(6.98\%\).
\end{itemize}

Retaining alternatives prevents the strongest current explanation from being
misrepresented as certainty.

\subsubsection{Withheld-event evaluation}

For each sealed project:
\begin{enumerate}
    \item a copied real source is revised;
    \item before and after SHA-256 values are recorded;
    \item AION predicts the affected artifacts before the transition is
    revealed;
    \item a confirmed or conflicted reconciliation event is withheld;
    \item the predicted state is compared with independent graph traversal;
    and
    \item a novel authority-revocation event tests model criticism.
\end{enumerate}

For a changed source \(s\) and dependency relation \(E\), the promoted model
predicts the transitive closure:
\[
I(s)
=
\left\{
v\;\middle|\;
\exists\text{ a directed path }s\xrightarrow{*}v\text{ in }E
\right\}.
\]

The cold control receives the same files and source-change signal but has no
learned relational model.  It invalidates every derived artifact.

\subsubsection{Novel-event criticism}

The sealed control event
\texttt{authority\_revoked\_without\_replacement} is outside the learned
event grammar.  AION returns:
\[
\texttt{EVENT\_OUTSIDE\_LEARNED\_MODEL\_FAMILY}
\]
and abstains.  It does not force the event into the nearest known transition.

\subsubsection{Phase 57 results}

\begin{table}[htbp]
\centering
\begin{tabular}{lr}
\hline
\textbf{Measure} & \textbf{Result}\\
\hline
Development projects & 5\\
Sealed projects & 4\\
Candidate models compared & 12\\
Competing explanations retained & 3\\
Withheld-event accuracy & 100\%\\
Weakest-family accuracy & 100\%\\
Novel-event criticism & 100\%\\
Unsafe novel-event forcing & 0\\
Unseen dependency topologies & 2\\
Provenance completeness & 100\%\\
Learned re-execution cost & 20\\
Cold re-execution cost & 26\\
Cost reduction versus cold & 23.08\%\\
Restart relearning & 0\\
\hline
\end{tabular}
\caption{Phase 57 sealed relational world-model results.}
\label{tab:phase57-results}
\end{table}

CAU promoted:
\[
\texttt{procedure\_richer\_world\_model\_0d316f0f69ad}.
\]

The complete HexCore history passed \(62/62\) tests.  The milestone was
committed as Git revision \texttt{c6e22199} and synchronized with remote
\texttt{main}.

\subsubsection{Persistent architecture}

The HexCore store now additionally contains:
\begin{itemize}
    \item \texttt{relational\_world\_models};
    \item \texttt{world\_model\_sessions};
    \item the selected theory and retained alternatives;
    \item posterior weights and development scores;
    \item dependency graphs;
    \item withheld predictions and revealed outcomes;
    \item model criticisms; and
    \item checksum-grounded event provenance.
\end{itemize}

Restart checks confirm that the learned model, every sealed session and the
promoted champion survive reconstruction with zero event relearning.

\subsubsection{Phase 57 boundary}

Phase 57 demonstrates transferable relational consequence learning over real
project files.  Source roles, dependency graphs, candidate operators,
revision injection and the correctness oracle remain engineered.  It is not
unrestricted world modelling, inference of arbitrary software semantics, or
causal discovery from raw physical reality.

\subsection{Combined interpretation}

The combined capability can be summarized as:
\[
\begin{aligned}
\text{begin a real-file project}
&\rightarrow \text{persist progress}\\
&\rightarrow \text{observe changed evidence}\\
&\rightarrow \text{predict affected consequences}\\
&\rightarrow \text{revise only dependent work}\\
&\rightarrow \text{abstain if explanation is inadequate}\\
&\rightarrow \text{verify against current sources}.
\end{aligned}
\]

Phase 56 establishes the operational control loop.  Phase 57 adds a learned,
transferable theory of project-state transitions.  Together they connect
knowledge ingestion, persistent project intelligence, world modelling,
metacognitive criticism and governed retention.

\subsection{Next stage: Phase 58 natural multimodal grounding}

Phase 58 should replace typed or synthetic modality contracts with natural
multimodal evidence:
\begin{itemize}
    \item real photographs and technical diagrams;
    \item charts, tables and spreadsheet regions;
    \item PDF page layout and visual annotations;
    \item software state and structured files;
    \item contradictory evidence across modalities; and
    \item temporal observations whose content changes.
\end{itemize}

The Phase 57 world model should predict cross-modal consequences, while exact
source evidence, independent execution, conflict handling and CAU remain the
final authority.
